\documentclass[11pt]{article}
\usepackage[margin=1in]{geometry}
\usepackage{amsmath,amssymb,amsfonts,mathtools,bm}
\usepackage[T1]{fontenc}
\usepackage[utf8]{inputenc}
\title{Inline and Display Mathematics}
\author{Source-backed grouped formula sample}
\date{}
\begin{document}
\maketitle
\section{Expressions}
\paragraph{Expression 1.} The following expression is taken from a source-backed formula corpus.

\begin{align*}\tilde{\psi}^{\dagger}_{\dot{\alpha}} \psi^{\dagger}_{\dot{\alpha}_{1}} \ldots \psi^{\dagger}_{\dot{\alpha}_{n}} \mid 0 \rangle = (-1)^{n} \psi^{\dagger}_{\dot{\alpha}_{1}} \ldots \psi^{\dagger}_{\dot{\alpha}_{n}}\psi^{\dagger}_{\dot{\alpha}} \mid 0 \rangle \vspace{-3pt}\end{align*}

\paragraph{Expression 2.} The following expression is taken from a source-backed formula corpus.

\begin{align*}H^F_D = \tilde{\lambda}^\mu(\tau) {\cal H}_\mu(\tau) - \frac 12 \tilde{\lambda}^{\mu \nu}(\tau) {\cal H}_{\mu \nu}(\tau) + \int{d^3\sigma \Big[ -A_\tau (\tau, \vec{\sigma}) \Gamma (\tau, \vec{\sigma}) + \mu_\tau (\tau,\vec{\sigma}) \pi^\tau (\tau, \vec{\sigma}) \Big] }, \end{align*}

\paragraph{Expression 3.} The following expression is taken from a source-backed formula corpus.

\begin{align*}X(\alpha_1,\alpha_2,\alpha'_1,\alpha'_2)=\exp \left[ i\alpha_1 q^1{\partial\over\partial q^1}+ i\alpha_2 q^2{\partial\over\partial q^2}+i\alpha'_1 \overline{q}^1{\partial\over\partial \overline{q}^1}+i\alpha'_2 \overline{q}^2{\partial\over\partial \overline{q}^2}\right]\end{align*}

\paragraph{Expression 4.} The following expression is taken from a source-backed formula corpus.

\begin{align*}\gamma = \gamma _{\underline{a}}^1I^{\underline{a}}+\gamma _{\underline{a} \underline{b}}^1I^{\underline{a}}I^{\underline{b}}+..., \text{ and } Q_\mu = q_{\mu ,\underline{a}}^1I^{\underline{a}}+q_{\mu ,\underline{a} \underline{b}}^2I^{\underline{a}}I^{\underline{b}}+...\end{align*}

\paragraph{Expression 5.} The following expression is taken from a source-backed formula corpus.

\begin{align*}\vec{B}=\frac{\Phi_0}{2\pi\lambda^2}\,\vec{b}\ \ ,\ \ \frac{\vec{E}}{c}=\frac{\Phi_0}{2\pi\lambda^2}\,\vec{e}\ \ ,\ \ F_{\mu\nu}=\frac{\Phi_0}{2\pi\lambda^2}\,f_{\mu\nu}\ \ ,\ \ A^0=\frac{\Phi_0}{2\pi\lambda}\,\varphi\ \ ,\ \ \vec{A}=\frac{\Phi_0}{2\pi\lambda}\,\vec{a},\end{align*}
\end{document}
